\chapter{Getting Started}
\label{ch:getting-started}

This chapter takes you from a fresh clone to an LED blinking under your own Ada,
in a few commands. The later chapters explain the runtime, the HAL and the
language techniques; this one just gets you running.

\section{What you need}

\begin{description}[style=nextline]
\item[An ESP32-S3 board.] Any devkit with the built-in USB-Serial-JTAG port (most
  of them). One USB cable is enough --- it carries the console \emph{and} the
  debugger.
\item[Alire (\texttt{alr}).] The Ada package manager fetches the toolchains the
  build needs: the \texttt{gnat\_xtensa\_esp32\_elf} cross-compiler, a native
  \texttt{gnat\_native} (for the host tools), and \texttt{gprbuild}. Install Alire,
  then make sure its toolchains are on your \texttt{PATH} for the shell you build
  from.
\item[Serial-port permission.] On Linux, add yourself to the \texttt{dialout}
  group (and install the udev rule for the USB-Serial-JTAG, see the debugging
  chapter) so you can open \texttt{/dev/ttyACM0} without \texttt{sudo}.
\end{description}

\section{Get the code}

The runtime depends on two git submodules (the \texttt{bb-runtimes} fork that
carries the ESP32-S3 bareboard runtime, and \texttt{xtensa-dynconfig}), so clone
\emph{recursively}:

\begin{lstlisting}[style=sh]
git clone --recurse-submodules \
    https://github.com/rowsail/ada_esp32s3.git
cd ada_esp32s3
#  (if you already cloned without submodules:)
git submodule update --init --recursive
\end{lstlisting}

\section{The \texttt{./x} dispatcher}

Everything goes through one script at the repository root. List the examples,
then build, flash and watch one:

\begin{lstlisting}[style=sh]
./x list                                   # the available examples + their profiles
./x run esp32s3_gpio0_blink -p /dev/ttyACM0   # build + flash + serial monitor
\end{lstlisting}

\texttt{./x run} builds the example, packages it (the Ada \texttt{esp\_elf2image}
host tool), writes it to the board (the Ada \texttt{esp\_flash} tool, over the
chip's ROM bootloader), resets, and streams the console. The first build also
generates the runtime for the chosen profile, so it takes a little longer; later
ones are incremental. The blink example toggles an LED and prints over the
console --- your first Ada running bare-metal on the S3.

The pieces are also separate commands when you want them:

\begin{lstlisting}[style=sh]
./x build   esp32s3_gpio0_blink            # cross-compile + link + package app.bin
./x flash   esp32s3_gpio0_blink -p /dev/ttyACM0
./x monitor -p /dev/ttyACM0                # just the serial console (115200)
./x clean   esp32s3_gpio0_blink
\end{lstlisting}

\section{Choosing a profile}

Most examples build under more than one runtime profile (Chapter~\ref{ch:profiles}).
Pass \texttt{--profile} to pick one; \texttt{auto} (the default) uses the
example's own:

\begin{lstlisting}[style=sh]
./x run esp32s3_heartbeat --profile embedded -p /dev/ttyACM0
\end{lstlisting}

\texttt{light-tasking} is the smallest, \texttt{embedded} adds exceptions, RAII
handles, a heap and the full HAL (the usual choice), and \texttt{full} adds the
complete Ada tasking model.

\section{Board configuration}

Each project owns a \texttt{config/board.ads} (flash and PSRAM sizes) that sizes
the image and the second-stage bootloader. Edit it directly, or:

\begin{lstlisting}[style=sh]
./x config esp32s3_gpio0_blink show
./x config esp32s3_gpio0_blink flash-size 4MB
\end{lstlisting}

\section{Debugging}

On-chip debugging needs the pinned OpenOCD + S3 GDB, fetched once; then
\texttt{./x debug} (or \texttt{F5} in VS Code) halts at your \texttt{Main}:

\begin{lstlisting}[style=sh]
./x get-debug-tools
./x debug esp32s3_gpio0_blink              # build, flash, OpenOCD + GDB, rest at Main
\end{lstlisting}

The Debugging chapter (Chapter~\ref{ch:debug}) covers GDB, decoding a Guru
Meditation, and the editor integration in full.

\section{Where to go next}

\begin{itemize}
\item The anatomy of an application and the runtime profiles ---
  Chapters~\ref{ch:anatomy}--\ref{ch:build}.
\item The peripheral HAL --- start at Chapter~\ref{ch:hal}; each driver has worked
  examples you can run with \texttt{./x run}.
\item The Ada techniques behind it all --- the Ada Language Techniques and
  Talking to Hardware chapters (Chapters~\ref{ch:embedded}
  and~\ref{ch:hardware}), including a Terminology section (\S\ref{sec:terms})
  for unfamiliar terms.
\end{itemize}

To start your \emph{own} project rather than an in-repo example, the
\texttt{esp32-ada} launcher scaffolds one in any empty folder (see the tooling
notes); it gives you the same \texttt{board.ads} and build flow without copying
the runtime.
